\section{Thread 状态机、等待队列与处理器现场}

前面的 v0.9--v1.1 路径用一个 PCB 同时保存页表根、文件状态、动态内核栈和
寄存器现场。这个结构不是错误：当一个进程恰有一条执行流时，它是最短、最容易
观察的实现，可以先证明 CPL3、CR3、TSS.RSP0、IRQ0 和 \code{iretq} 真的
连成闭环。问题出现在继续演进时。若两个执行流共享同一个地址空间，页表根应当
只有一份，但 RIP、RSP、内核栈和浮点寄存器必须各有一份；一个控制块已经无法
同时表达两种所有权。

\begin{keyidea}
Process 回答“哪些资源被一组执行流共享”；Thread 回答“哪一份现场有资格取得
CPU”。调度器只选择 Thread，地址空间、描述符和退出关系则由 Process 统一
拥有。
\end{keyidea}

这种拆分不是为了模仿某个接口名称，而是由寿命差异强制得到。Thread 可以在
同一 Process 内独立退出，其他 Thread 继续使用原地址空间；只有最后一个
live Thread 离开，Process 才能发布退出结果并进入 Zombie。反过来，地址
空间也不能在仍有 Thread 可能返回用户态时销毁。

\begin{pdfdiagram}
  \node[os state] (process) {Process\\PID / AddressSpace\\描述符与退出结果};
  \node[os block,below left=of process] (threada) {Thread A\\TID / 帧 / 栈\\FXSAVE};
  \node[os block,below right=of process] (threadb) {Thread B\\TID / 帧 / 栈\\FXSAVE};
  \node[os hardware,below=of threada] (cpu) {单个 BSP\\当前 CPU};
  \node[os state,below=of threadb] (wait) {run queue\\WaitQueue};
  \draw[os arrow] (process) -- node[left]{owns} (threada);
  \draw[os arrow] (process) -- node[right]{owns} (threadb);
  \draw[os arrow] (threada) -- (cpu);
  \draw[os arrow] (threadb) -- (wait);
  \draw[os arrow] (wait) -- (cpu);
\end{pdfdiagram}
\diagramnote{Process 不进入 run queue；Thread 不独占共享地址空间。当前只有一个 BSP，但对象边界已经不依赖“四个程序”这一偶然数量。}

\topic{两级对象分别保存什么}\par
\begin{osgridlongtable}{L{0.20\textwidth}L{0.31\textwidth}L{0.37\textwidth}}
对象 & 当前字段组 & 所有权不变量\\ \hline
Process &
\code{ProcessId}、Alive/Zombie、页表根、Thread 链、live/exited 数 &
PID 不等于槽位；页表根非零；Thread 链计数与三类计数一致\\ \hline
Thread &
\code{ThreadId}、所属 Process、调度状态、内核栈槽、用户栈、TLS、signal mask &
TID 独立单调；每个活动 Thread 恰属于一个 Process\\ \hline
Thread 现场 &
176 B 通用/IRET 帧、512 B FXSAVE 区 &
帧完整位于本 Thread 动态内核栈；FXSAVE 区 16 B 对齐\\ \hline
Thread 队列 &
Process 链、双向 run queue 链、WaitQueue 链 &
一个 Thread 不能同时是 Running、Ready 和 waiter\\ \hline
统计 &
run tick、dispatch、block、wake &
当前量与累计 create/discard/reap 分开\\ \hline
\end{osgridlongtable}

PID 与 TID 都用显式 \code{uint64_t} 包装，零值表示无身份；
\code{UINT64_MAX} 只表示无槽索引。身份和存储位置使用不同哨兵，可以避免把
“数组末端”意外解释为一个合法 PID。槽位在 reap 后会复用，身份计数器不回退，
所以旧 TID 不会因为新 Thread 恰好落入同一槽而重新有效。

当前固定数组是 freestanding 早期阶段的存储后端，不是公开语义。容量档可以
把有效前缀限制为 4、128 或 512 个 Thread；调度状态机和身份分配仍是同一份
代码。以后把槽替换为 type cache 时，PID/TID、队列和退出协议不需要改变。

\topic{两个状态机比一个巨型枚举更能表达寿命}
Thread 状态机为：

\[
  \text{Unused}\rightarrow\text{Ready}
  \leftrightarrow\text{Running}
  \rightarrow\text{Blocked}
  \rightarrow\text{Ready},
\]
\[
  \text{Running}\rightarrow\text{Exited}\rightarrow\text{Unused}.
\]

Process 状态机为：

\[
  \text{Unused}\rightarrow\text{Alive}
  \xrightarrow{\mathrm{live\ threads}=0}
  \text{Zombie}\rightarrow\text{Unused}.
\]

\begin{pdfdiagram}
  \node[os state] (unused) {Thread\\Unused};
  \node[os block,right=of unused] (ready) {Ready};
  \node[os hardware,right=of ready] (running) {Running};
  \node[os state,below=of running] (blocked) {Blocked};
  \node[os block,right=of running] (exited) {Exited};
  \draw[os arrow] (unused) -- node[above]{create} (ready);
  \draw[os arrow] (ready) -- node[above]{dispatch} (running);
  \draw[os arrow] (running) to[bend left] node[above]{quantum/yield} (ready);
  \draw[os arrow] (running) -- node[right]{wait} (blocked);
  \draw[os arrow] (blocked) -- node[below]{wake winner} (ready);
  \draw[os arrow] (running) -- node[above]{exit} (exited);
  \draw[os arrow] (exited) to[bend left] node[below]{safe reap} (unused);
\end{pdfdiagram}
\diagramnote{Exited 表示不再执行，却不表示栈已经可以在当前 RSP 上销毁；状态发布与资源安全点是两个时刻。}

每次完整校验都重新扫描对象，而不是相信增量计数。核心守恒式为：

\[
\begin{aligned}
T_{\mathrm{owned}} &=
T_{\mathrm{ready}}+T_{\mathrm{running}}+
T_{\mathrm{blocked}}+T_{\mathrm{exited}},\\
P_{\mathrm{owned}} &=P_{\mathrm{alive}}+P_{\mathrm{zombie}},\\
P_i.\mathrm{thread\_count} &=
P_i.\mathrm{live\_thread\_count}+
P_i.\mathrm{exited\_thread\_count},\\
T_{\mathrm{running}} &\le 1.
\end{aligned}
\]

若 Process 为 Zombie，它的 live 数必须为零；若 Thread 为 Running，它必须
等于全局 current 索引，且不得仍带 run queue 或 WaitQueue 链接。Ready
Thread 必须恰好在 run queue 出现一次。扫描还限制遍历步数不超过 Thread
容量，使环链被诊断为损坏，而不是让校验函数自身无限循环。

\topic{为什么 Thread 与 Process 要分两次 reap}
当前正常用户程序每个 Process 只有一个初始 Thread，但回收仍执行完整两级
协议：

\begin{enumerate}[label=\arabic*.]
  \item 系统调用或用户异常先把当前 Thread 发布为 Exited；
  \item 最后一个 live Thread 使所属 Process 进入 Zombie；
  \item 当前 CPU 切回永久内核 CR3，Process 地址空间不再可能被返回使用；
  \item 汇编恢复 \code{ExecuteProcesses} 保存的永久内核栈；
  \item 安全点确认真实 RSP 不在待销毁栈内，销毁 Thread 动态栈；
  \item \code{ReapExitedThread} 从 Process 所有权链移除 Thread；
  \item 只有 \code{thread\_count=0} 才能 \code{ReapZombieProcess}。
\end{enumerate}

若先 reap Process，Thread 的所属对象会成为悬空索引；若在退出处理函数中直接
撤销当前内核栈的 PTE，函数的下一次压栈或返回就会页故障。两级状态与安全点
共同把“逻辑停止”“硬件不再使用”和“存储已归还”分开。

\topic{侵入式队列让调度热路径不依赖分配成功}
Ready 队列把前后槽索引直接放进 Thread。入队写
\code{previous\_run\_thread\_index} 与
\code{next\_run\_thread\_index}，维护 head/tail；出队清除两个链接后才
把 Thread 变成 Running。Process 的 Thread 链只需单向遍历，WaitQueue
也用单向 FIFO。这样没有额外 node 对象，也不会在 IRQ 或退出边界调用 heap。

\begin{contractbox}
队列成员链接是状态的一部分。一个 Ready Thread 的双向链接必须与全局
head/tail 相符；Blocked Thread 的 run 链必须为空；Exited Thread 的 run
与 wait 链都必须为空。修改状态和链接必须在同一个调度提交区完成。
\end{contractbox}

round-robin 的时间片仍为四个 PIT tick。IRQ0 到来时先记录当前 Thread 的
run tick；预算未到或没有其他 Ready 时返回原帧。预算到且队列非空时：

\begin{enumerate}[label=\arabic*.]
  \item 保存当前 Thread 的 176 B 帧指针和 512 B 扩展现场；
  \item 把当前 Thread 置 Ready 并追加到队尾；
  \item 弹出队首，置为 Running，递增 dispatch/preemption；
  \item 激活目标 Process 的 CR3；
  \item 把目标 Thread 动态栈顶写入 TSS.RSP0；
  \item 恢复目标 Thread 的扩展现场；
  \item 汇编从目标 176 B 帧逆序弹寄存器并执行 \code{iretq}。
\end{enumerate}

\begin{pdfdiagram}
  \node[os hardware] (irq) {PIT IRQ0\\CPL3 入栈};
  \node[os block,right=of irq] (save) {保存通用帧指针\\FXSAVE64 current};
  \node[os state,right=of save] (decision) {ThreadScheduler\\队尾入队/队首出队};
  \node[os block,below=of decision] (activate) {CR3(Process)\\TSS.RSP0(Thread)};
  \node[os hardware,left=of activate] (restore) {FXRSTOR64 next\\IRETQ};
  \draw[os arrow] (irq) -- (save);
  \draw[os arrow] (save) -- (decision);
  \draw[os arrow] (decision) -- (activate);
  \draw[os arrow] (activate) -- (restore);
\end{pdfdiagram}
\diagramnote{调度器只返回索引和决策；CR3、TSS、FXSAVE 与帧地址验证留在目标机 ProcessRuntime。}

\topic{WaitQueue 解决的是竞争完成，不只是“睡一下”}
一个读者等待数据时，至少有五种事件可能结束等待：条件满足、deadline 到达、
signal、对象关闭、映射取消。若五条路径各自执行一次“设为 Ready”，同一
Thread 可能被插入 run queue 多次。链表随后不是简单地“多运行一次”，而是
可能形成环、覆盖 tail 或让统计永久失衡。

v1.2 为每个等待对象建立 \code{WaitQueueId}、head/tail、当前 waiter 数、
累计 enqueue/wake、close 次数和 closed 状态。阻塞事务为：

\begin{lstlisting}[language=C++,caption={等待提交的语义顺序}]
Running
  -> wait_condition = requested_condition
  -> wake_reason = None
  -> wait_queue = &object_queue
  -> append FIFO waiter
  -> Blocked
  -> select next Ready Thread
\end{lstlisting}

唤醒事务只接受以下前置状态：

\[
  \mathrm{state=Blocked}\land
  \mathrm{queue=target}\land
  \mathrm{wake\_reason=None}.
\]

第一个完成者从 WaitQueue 移除 Thread，写入具体 WakeReason，再追加到 Ready
队尾。第二个完成者发现前置状态已不成立，得到
\code{WakeAlreadyResolved}；它不能再次入队。这里的“单赢家”不是指只有
一种事件能发生，而是只有一种事件能取得状态转换的提交权。

\begin{osgridlongtable}{L{0.27\textwidth}L{0.25\textwidth}L{0.36\textwidth}}
WakeReason & 典型来源 & 调用者恢复后的含义\\ \hline
ConditionSatisfied & 数据到达、空间释放、Mutex handoff & 重新检查条件并尝试操作\\ \hline
Timeout & 单调 deadline 到达 & 无进度时返回超时；本阶段仅冻结原因\\ \hline
Signal & 可中断等待收到信号 & 按系统调用重启规则处理；后续阶段实现\\ \hline
ObjectClosed & pipe/descriptor/队列关闭 & 重新检查 EOF、broken pipe 或关闭错误\\ \hline
Cancelled & futex 地址 unmap 等取消 & 资源身份已失效，不得继续使用旧 key\\ \hline
\end{osgridlongtable}

关闭不是简单地设置布尔值。实现先令 queue closed，使新阻塞立即失败；再按
FIFO 以 ObjectClosed 唤醒所有已登记 waiter。若先逐个唤醒、最后才 closed，
中间可能重新进入一个新 waiter，关闭操作就无法形成明确边界。

\topic{检查条件与入队为什么必须不可分割}
经典丢失唤醒序列是：

\begin{enumerate}[label=\arabic*.]
  \item 读者观察“无数据”；
  \item 生产者写入数据并尝试唤醒，但队列尚为空；
  \item 读者随后入队并睡眠；
  \item 再无事件发生，数据存在而读者永久睡眠。
\end{enumerate}

当前单 BSP 的 \code{INT 0x80} interrupt gate 会清 IF，调度提交再由
\code{IrqSaveSpinLock} 包围。Try/Wait 用户包装在 Wait 系统调用中重新检查
条件，只有仍不满足才入队。SMP 后关本地中断不能阻止另一 CPU，必须让对象
条件锁与 WaitQueue 提交遵守同一锁序；当前单赢家状态机可以原样保留。

\topic{三类锁服务三种不同等待成本}\par
\begin{osgridlongtable}{L{0.22\textwidth}L{0.25\textwidth}L{0.18\textwidth}L{0.23\textwidth}}
原语 & 使用位置 & 能否睡眠 & v1.2 关键语义\\ \hline
SpinLock & Thread 上下文短提交区 & 否 & exchange-acquire、store-release、内层 PAUSE\\ \hline
IrqSaveSpinLock & 与本 BSP IRQ 共享状态 & 否 & 保存 IF、CLI、取锁、释放后恢复原 IF\\ \hline
Mutex & 可能较长的 Thread 临界区 & 是 & 竞争者进 WaitQueue，解锁直接 handoff\\ \hline
\end{osgridlongtable}

SpinLock 的 acquire 防止临界区读取被移到取锁前，release 防止临界区写被移到
解锁后。\code{pause} 只减少忙等对执行资源的压力，不提供互斥或内存屏障。
IrqSaveSpinLock 不能在释放时无条件 \code{sti}：调用者进入前可能已经关中断，
所以必须恢复保存的 IF。

Mutex 竞争时不能持有任何 spinlock。否则当前 Thread 睡眠后，另一个需要该
spinlock 才能满足条件的 Thread 永远无法前进。v1.2 的
\code{CurrentSpinLockDepth} 是单 BSP 防误用门禁；深度非零时 Mutex 返回
\code{SpinLockHeld}，不调用调度器。

解锁存在 waiter 时使用直接交接：队首 Thread 在变 Ready 的同时成为逻辑
owner，并带有 handoff pending。它实际运行后用一次 \code{TryLock} 确认
handoff。若先把 owner 清空再普通唤醒，另一 Ready Thread 可能在原队首运行
之前抢到 Mutex，形成不必要的插队和饥饿。

\topic{x87 到 SSE2 的历史解释了为什么现场不只是一组通用寄存器}
8086 最初没有片上浮点单元，8087 是独立数值协处理器。它采用八槽寄存器栈，
并拥有 control、status、tag、instruction pointer 与 data pointer。后来的
x86 把 FPU 集成进处理器，却为了兼容保留了这套可见状态。MMX 又复用 x87
物理寄存器，所以 MMX 与 x87 也不能被当作互不相关的临时值。

SSE 在 Pentium III 时代加入 128 位 XMM 与 MXCSR；SSE2 扩展整数和双精度
操作，并成为 x86-64 软件实际依赖的基线之一。64 位模式可见 XMM0--XMM15，
而不是 32 位时代常见的八个。CPU 切换 CR3、RSP 或执行 \code{iretq} 时都
不会自动保存这些状态。

\begin{failurebox}
若 Thread A 在 XMM0 中留下密钥或计算中间值，随后只保存通用寄存器就切到
Thread B，B 可以直接读到 A 的 XMM0。这既会产生错误结果，也会形成跨安全域
信息泄漏。禁止普通 C++ 生成 SSE 只能缩小测试噪声，不能代替内核保存协议。
\end{failurebox}

\topic{CPUID 能力与 OS 控制位必须同时成立}
\code{CPUID.01H:EDX} 的三位是本阶段硬门槛：

\begin{osgridlongtable}{L{0.12\textwidth}L{0.17\textwidth}L{0.59\textwidth}}
位 & 名称 & 意义\\ \hline
24 & FXSR & 处理器支持 FXSAVE/FXRSTOR\\ \hline
25 & SSE & 存在 XMM 与 MXCSR 状态\\ \hline
26 & SSE2 & 满足当前 64 位用户程序的扩展指令基线\\ \hline
\end{osgridlongtable}

能力存在仍不代表 OS 已准备好。v1.2 配置：

\begin{osgridlongtable}{L{0.14\textwidth}L{0.24\textwidth}L{0.14\textwidth}L{0.36\textwidth}}
寄存器 & 位 & 当前值 & 原因\\ \hline
CR0 & MP(1) & 1 & 让 WAIT/FWAIT 与 TS 规则正确配合\\ \hline
CR0 & EM(2) & 0 & 不使用软件仿真，允许真实 x87 指令\\ \hline
CR0 & TS(3) & 0 & 不使用基于 \#NM 的惰性切换\\ \hline
CR0 & NE(5) & 1 & x87 错误经原生异常，不走旧式 IRQ13\\ \hline
CR4 & OSFXSR(9) & 1 & 声明 OS 管理 FXSAVE 现场\\ \hline
CR4 & OSXMMEXCPT(10) & 1 & SIMD 浮点异常经 \#XM\\ \hline
CR4 & OSXSAVE(18) & 0 & 禁止尚未保存的 XSAVE/AVX 状态\\ \hline
\end{osgridlongtable}

AVX 的 YMM 上半部不在 FXSAVE 中。若开启 OSXSAVE 却不读取 CPUID leaf 0xD、
不设置 XCR0，也不按动态大小分配区域，内核就会出现部分支持。因此本阶段不是
“忘了做 AVX”，而是主动把不可维护的状态面关闭。

\topic{512 字节 FXSAVE 区是一份硬件规定的二进制结构}
\code{FXSAVE64} 要求内存操作数按 16 字节对齐。v1.2 用
\code{alignas(16)} 定义恰好 512 字节的 \code{FxSaveArea}，并用两个
\code{static_assert} 锁定大小和对齐。主要布局为：

\begin{osgridlongtable}{L{0.14\textwidth}L{0.15\textwidth}L{0.48\textwidth}L{0.13\textwidth}}
偏移 & 长度 & 内容 & 验收关注\\ \hline
\code{0x000} & 8 B & FCW、FSW、abridged FTW、FOP & x87 控制/状态\\ \hline
\code{0x008} & 16 B & 64 位 FIP 与 FDP & 指令/数据指针\\ \hline
\code{0x018} & 8 B & MXCSR 与 MXCSR\_MASK & SIMD 控制/能力\\ \hline
\code{0x020} & 128 B & ST0--ST7/MMX，每槽 16 B & x87 数据\\ \hline
\code{0x0A0} & 256 B & XMM0--XMM15，每槽 16 B & SSE2 数据\\ \hline
\code{0x1A0} & 96 B & 保留或软件可用尾区 & 总长补足 512 B\\ \hline
\end{osgridlongtable}

x87 每个寄存器的有效值为 80 位，但槽宽是 16 字节；不能把这 128 字节直接
解释为八个 \code{double}。FXSAVE 使用 abridged tag word，也不能与老式
FSAVE 的完整 tag word 混用。当前代码不手工解析字段，只把整块交给硬件，
避免软件复制遗漏保留状态。

初始化先执行 \code{fninit}，再载入体系结构常见默认
\code{MXCSR=0x1F80}，最后 FXSAVE64 到全局初始模板。每次创建 Thread 只
复制模板，不重复改变当前 CPU 浮点环境。保存与恢复接口先检查初始化状态和
区域对齐，成功后才增加冷路径累计计数。

\topic{为什么选择 eager 而不是 lazy FPU switching}
历史内核常在上下文切换时设置 CR0.TS。新任务第一次执行浮点指令触发
\#NM，内核此时才保存旧 FPU owner 并恢复新 owner。浮点使用稀少、保存指令
昂贵时，这种 lazy 策略能节省工作。

代价是新增一个与 scheduler current 不同的“当前 FPU owner”，还要处理
\#NM 嵌套、内核代码是否可用浮点、退出时 owner 清理，以及推测执行时代的
敏感状态边界。本项目选择 eager：抢占、阻塞和退出保存 current，激活目标
Thread 恢复 next。多做一次固定 512 字节硬件事务的成本，换来没有隐藏 owner
的直接状态机。

\topic{四组用户模式让保存协议接受真实硬件审问}
布局单元测试只能证明结构大小，累计计数只能证明函数被调用；二者都不能证明
保存区属于正确 Thread。v1.2 的四个 Ring 3 程序按 PID 安装互不相同的模式：

\begin{osgridlongtable}{L{0.19\textwidth}L{0.27\textwidth}L{0.42\textwidth}}
状态 & 模式 & 覆盖目的\\ \hline
XMM0 & 每 PID 两个不同 64 位分量 & 覆盖常用低 XMM 数据\\ \hline
XMM15 & 另一组两个 64 位分量 & 覆盖 64 位模式高 XMM 区\\ \hline
MXCSR & 四个有效舍入/控制组合 & 证明控制状态也隔离\\ \hline
x87 control word & 四个不同舍入控制字 & 覆盖 x87 控制字段\\ \hline
ST0 & 1.0、2.0、3.0、4.0 & 覆盖 x87 数据栈且可精确比较\\ \hline
\end{osgridlongtable}

普通用户 C++ 继续使用 \code{-mno-sse -mno-sse2}，只有
\filepath{source/user/src/extended_state.asm} 修改被观察状态。安装函数
立即自检，排除表索引错误；完成函数在写日志之前校验一次，输出
\code{EXTENDED_STATE_ISOLATED} 后再校验一次，证明系统调用本身没有污染。

四个程序经过不同调度路径：

\begin{itemize}
  \item Shell 等待真实 PS/2 键盘字符，经历无 Ready 时的 idle 与唤醒；
  \item producer 在 64 字节管道写满时阻塞；
  \item consumer 在管道为空时阻塞并在写端关闭后观察 EOF；
  \item worker 执行有界计算，由真实 1000 Hz PIT 多次抢占。
\end{itemize}

QEMU 必须精确看到四次隔离标记，并看到 Kernel save/restore 均非零。再使用
\code{qemu64,-sse2} 运行同一磁盘；此时必须只到达
\code{EXTENDED_STATE_UNSUPPORTED}，不能进入 GDT、用户态或 READY。
成功与能力失败共享同一实现，避免“测试专用内核”掩盖真实边界。

\topic{三档容量事务证明规格不是常量表}\par
\begin{osgridlongtable}{L{0.20\textwidth}L{0.13\textwidth}L{0.16\textwidth}L{0.15\textwidth}L{0.20\textwidth}}
档位 & RAM & Process & Thread & 单 Process Thread\\ \hline
bootstrap & 64 MiB & 4 & 4 & 1\\ \hline
functional & 256 MiB & 64 & 128 & 32\\ \hline
capacity & 64 GiB & 256 & 512 & 64\\ \hline
\end{osgridlongtable}

只初始化 512 个空结构体不能证明系统能承载 512 个 Thread。启动容量事务为
每个 Process 创建真实用户页表根，为每个 Thread 创建六页 KVA 动态栈：
上下 guard 保持 not-present，中间四页各有 order-0 物理后备；同时初始化
每份 FXSAVE 区。至少一个 Process 实际达到当前档位的单进程 Thread 上限。

创建后的中间快照必须观察：

\[
\begin{aligned}
\mathrm{process\_count} &= P_{\mathrm{limit}},\\
\mathrm{thread\_count} &= T_{\mathrm{limit}},\\
\mathrm{active\_stack\_count} &= T_{\mathrm{limit}},\\
\mathrm{mapped\_stack\_pages} &= 4T_{\mathrm{limit}},\\
\mathrm{guard\_pages} &= 2T_{\mathrm{limit}}.
\end{aligned}
\]

全部 Thread 随后经调度器进入 Exited，Process 进入 Zombie。清理先销毁栈并
立即清除回滚活动标志，再 reap Thread；再销毁页表根并立即清除根记录，最后
reap Process。立即清标志很关键：若后续 scheduler reap 失败，统一清理函数
不能再次销毁已经归还的栈或页表。

\begin{failurebox}
页表根创建成功后必须先登记到回滚数组，再调用 CreateProcess。若顺序相反，
调度器拒绝会让一个尚未进入 Process 槽的真实根失去所有权记录。容量越大，
这种“提交之前的资源”越容易在罕见失败分支中泄漏。
\end{failurebox}

前后 \code{ResourceSnapshot} 比较 frame、buddy、heap、KVA、动态栈和补充
Process/Thread 当前量，变化掩码必须为零。历史累计允许增加；当前所有权必须
恢复。为使 512 个栈不会先耗尽次级元数据，KVA 描述符容量提升为 1024，
内核栈槽容量提升为 512。

\topic{测试不是数量装饰，而是四个不同观察位置}\par
\begin{osgridlongtable}{L{0.18\textwidth}L{0.31\textwidth}L{0.39\textwidth}}
层次 & 输入 & 能证明的性质\\ \hline
单元 & 非法容量、256/512 对象、五种 WakeReason、Mutex/锁 &
失败原子性、状态机、FIFO、身份与计数\\ \hline
集成 & 三 Process、六 Thread、48 tick、阻塞/唤醒/退出 &
多 Thread 公平性与两级所有权组合\\ \hline
随机 & 固定种子 100000 步操作序列 &
长序列槽复用、三条 WaitQueue、累计与当前量对照\\ \hline
QEMU & 真实 CPUID、CR0/CR4、PIT、PS/2、管道和 Ring 3 &
硬件现场隔离、三档容量、失败停止与资源回收\\ \hline
\end{osgridlongtable}

随机模型的种子为 \code{0x5448524541445632}。每一步从 create process、
create thread、start/tick/yield、block、wake、exit、reap 和 discard 中
选择一个动作，再与独立参考状态比较。测试不是在 100000 步结束时才
\code{Validate}；每一步都检查 current Thread、三条 WaitQueue、每个状态
集合、PID/TID 单调性和累计计数，因此首次破坏位置可以稳定复现。

完整回归为 99 项 CTest。新增扩展现场布局单元测试和 SSE2 缺失系统用例；
旧 ProcessScheduler 的三项测试被等层替换为 ThreadScheduler 单元、集成和
十万步随机测试，而全部历史启动、损坏镜像、用户异常、文件持久化、命名和
教材门禁继续存在。

\topic{按目录走读 v1.2 实现}\par
\begin{osgridlongtable}{L{0.43\textwidth}L{0.45\textwidth}}
文件 & 阅读问题\\ \hline
\filepath{process/thread_scheduler.hpp/.cpp} &
Process/Thread 状态、侵入式链和 Validate 如何彼此约束？\\ \hline
\filepath{process/wait_queue.hpp/.cpp} &
为什么队列只暴露生命周期与统计，实际链接由 scheduler 维护？\\ \hline
\filepath{sync/spin_lock.hpp/.cpp} &
acquire/release、PAUSE、IF 保存和持锁深度在哪里建立？\\ \hline
\filepath{sync/mutex.hpp/.cpp} &
直接 handoff 如何避免 owner 空窗与插队？\\ \hline
\filepath{arch/extended_state_layout.hpp/.cpp} &
哪些能力和布局可在宿主测试而不执行特权指令？\\ \hline
\filepath{arch/extended_state.hpp/.cpp} &
CPUID、CR0/CR4、初始模板和累计计数如何组合？\\ \hline
\filepath{arch/architecture.asm} &
FXSAVE64/FXRSTOR64 为什么保持为三个极小 Intel 语法边界？\\ \hline
\filepath{process/process_runtime.cpp} &
纯调度决定如何落实为 CR3、TSS.RSP0、动态栈和扩展现场？\\ \hline
\filepath{user/src/extended_state.asm} &
四组 XMM/MXCSR/x87 模式如何在不依赖 C++ 浮点代码的情况下验证？\\ \hline
\end{osgridlongtable}

推荐先读 layout 与单元测试，再读 scheduler 和随机模型，最后进入
ProcessRuntime 与两份汇编。直接从目标运行时开始，会同时面对页表、栈、
队列和 FXSAVE，难以判断一个分支属于策略还是硬件落实。

\topic{本阶段明确没有做什么}
v1.2 的正常用户进程仍只有一个初始 Thread；用户
\code{CreateThread/Join}、TLS ABI 和 private futex 在后续阶段开放。当前
Process/Thread 存储仍有 256/512 编译期上界，但身份、状态和等待协议不依赖
槽位。内核仍是单 BSP、内核不可抢占模型，没有 per-CPU run queue、远端 TLB
shootdown 或 SMP 锁竞争。

\code{INT 0x80} 仍是唯一用户入口。下一阶段先建立 CpuLocal、可信入口栈、
\code{swapgs} 与统一 UserContext，再加入 \code{syscall/sysret}；不能在
Thread 模型尚未冻结时同时改写架构入口。AVX/XSAVE 也保持关闭。清晰记录未做
的边界，使后续阶段知道哪些不变量可以复用，哪些能力仍必须从失败路径开始设计。

\begin{keyidea}
v1.2 的完成不是“能同时放 512 个结构体”。它同时证明：Thread 是唯一调度
实体，等待只有一个完成赢家，x87/SSE2 不跨执行流泄漏，三档真实资源能够
完整取得并归还，且这些事实可以由宿主模型和从复位启动的 QEMU 整机共同重现。
\end{keyidea}
